\chapter{Difficulty Reaching Female Orgasm}

\section*{Definition and Overview}
Difficulty reaching orgasm, also called anorgasmia, is the persistent or recurrent difficulty achieving orgasm despite adequate stimulation and arousal, and it can cause distress or affect relationships. It is one of the most common sexual difficulties in women, and it can occur in all situations or only in specific ones, such as with a partner but not during masturbation. Many factors contribute, including psychological, relational, and physical ones, and most cases respond well to education, communication, and specific techniques. This chapter explains the causes and treatment of difficulty reaching orgasm in women.

\section*{Epidemiology}
Difficulty reaching orgasm is common: around 10--40\% of women report some difficulty reaching orgasm, and about 5--10\% rarely or never do. It is more common with age and after menopause, and in women with depression, anxiety, or relationship problems. Many women do not reach orgasm from penetrative sex alone and need direct clitoral stimulation. Because it is often not discussed, many women suffer in silence without seeking help, despite effective treatments being available.

\section*{Aetiology and Pathophysiology}
\begin{itemize}
    \item \textbf{Insufficient stimulation:} many women need direct, prolonged clitoral stimulation
    \item \textbf{Psychological factors:} anxiety, stress, performance pressure, and body image concerns
    \item \textbf{Relationship factors:} poor communication, conflict, and lack of emotional intimacy
    \item \textbf{Depression and anxiety:} mental health conditions and their treatments affect sexual function
    \item \textbf{Medicines:} antidepressants, especially SSRIs, can delay or prevent orgasm
    \item \textbf{Medical conditions:} diabetes, neurological conditions, pelvic surgery, and hormonal changes
    \item \textbf{Alcohol and drugs:} can impair arousal and orgasm
\end{itemize}

\section*{Clinical Features}
\begin{itemize}
    \item Orgasm takes a long time or does not happen despite arousal
    \item Primary anorgasmia: never having had an orgasm
    \item Secondary anorgasmia: previously able to orgasm, now having difficulty
    \item Situational anorgasmia: orgasm possible with masturbation but not with a partner, or in some situations but not others
    \item Distress, frustration, and reduced sexual satisfaction
    \item Avoidance of sexual activity and relationship strain
\end{itemize}

\section*{Differential Diagnosis}
\begin{itemize}
    \item Low sexual desire, which is a different problem from difficulty reaching orgasm
    \item Vaginismus and painful intercourse, which can coexist and interfere with arousal
    \item Depression, anxiety, and medication effects
    \item Hormonal changes, including menopause and postpartum changes
    \item Medical conditions affecting pelvic nerves and blood flow
\end{itemize}

\section*{When to Seek Medical Care}
Women who are distressed by difficulty reaching orgasm should discuss it with a doctor, who can review medications, check for underlying conditions, and refer to a sexual health physician, gynaecologist, or psychologist.

\textbf{Call 000 (or local emergency services) for:}
\begin{itemize}
    \item Sudden loss of sexual function with pelvic pain, weakness, or numbness, which may need urgent assessment
    \item Any medical emergency
\end{itemize}

\section*{Diagnosis and Investigations}
\begin{itemize}
    \item \textbf{Detailed history:} when the problem occurs, arousal and stimulation patterns, relationship context, and medicines
    \item \textbf{Psychological assessment:} mood, stress, anxiety, and body image
    \item \textbf{Physical examination:} pelvic examination where relevant
    \item \textbf{Blood tests:} hormone levels, thyroid function, and blood glucose where indicated
    \item \textbf{Medication review:} identifying drugs that affect orgasm
\end{itemize}

\section*{Management}
\begin{itemize}
    \item \textbf{Education:} understanding that direct clitoral stimulation is normal and necessary for most women
    \item \textbf{Communication:} discussing needs and preferences with a partner
    \item \textbf{Sensate focus:} structured exercises to reduce pressure and build intimacy
    \item \textbf{Masturbation practice:} exploring what stimulation works best
    \item \textbf{Vibrators and aids:} can help many women reach orgasm
    \item \textbf{Psychological therapy:} cognitive-behavioural therapy and sex therapy for anxiety and relationship issues
    \item \textbf{Medication review:} adjusting antidepressants under medical supervision where they are the cause
    \item \textbf{Managing underlying conditions:} treatment of hormonal, neurological, or other medical problems
\end{itemize}

\section*{Complications}
\begin{itemize}
    \item Distress, low self-esteem, and reduced sexual satisfaction
    \item Relationship conflict and avoidance of intimacy
    \item Depression and anxiety
    \item Difficulty conceiving in some cases, when sex is avoided
\end{itemize}

\section*{Prognosis and Outcomes}
Most women with difficulty reaching orgasm improve significantly with education, communication, and specific techniques, and the prognosis is good. Situational difficulty responds especially well to learning what stimulation works and sharing that with a partner. Where medicines or medical conditions are involved, addressing the cause often restores function. With professional support and open communication, most women and couples achieve satisfying sexual lives.

\section*{Prevention and Self-Care}
\begin{itemize}
    \item Learn what stimulation works for you, including masturbation
    \item Communicate openly with your partner about what you need
    \item Reduce performance pressure and focus on pleasure, not goals
    \item Manage stress and prioritise time for intimacy
    \item Discuss possible sexual side effects before starting medicines
    \item Limit alcohol and avoid recreational drugs
    \item Seek help early if it is causing distress
\end{itemize}

\section*{Sources and Further Reading}
The following reputable sources provide further information for healthcare students and patients seeking reliable, current advice about difficulty reaching orgasm.

\begin{itemize}
    \item Healthdirect Australia. \textit{Female sexual problems: information and support.} https://www.healthdirect.gov.au/
    \item Jean Hailes for Women's Health. \textit{Sexual health and orgasm.} https://www.jeanhailes.org.au/
    \item NHS. \textit{Female sexual problems and orgasm difficulty.} https://www.nhs.uk/conditions/sexual-problems-in-women/
    \item Mayo Clinic. \textit{Female sexual dysfunction: causes and treatment.} https://www.mayoclinic.org/
    \item MSD Manual. \textit{Female orgasmic disorder: professional overview.} https://www.msdmanuals.com/
    \item MedlinePlus. \textit{Orgasmic dysfunction.} https://medlineplus.gov/
\end{itemize}
